%% toc: P3 39 39.2
\chapter{Definitional Complexity and Finite-Variable Logics}
\label{sec:p03-w03-p03-055-definitional-complexity}

A property of finite structures can be classified by the logics able to
define it and, inside one logic, by how much of the logic a definition
needs. This chapter measures that need in two ways: by quantifier rank, the
depth of quantifier nesting, and by the number of distinct variables a
formula uses. Each measure has a game that decides when two structures cannot
be told apart within the measure: the Ehrenfeucht--Fra\"iss\'e game for
quantifier rank and the pebble game for the number of variables. The
rank game is proved in \cref{sec:p03-w03-p03-038-comparison-games}; this
chapter adds the converse that makes the rank game a complete method,
proves the pebble-game characterization of the finite-variable logics
and its refinement in which rounds and pebbles are both bounded,
derives the corresponding definability criteria, shows that both
hierarchies are strict, and works through cardinality, parity,
connectivity and 2-colorability as test cases, ending with a procedure
for classifying a concept and a table of the cases treated.

\emph{Inputs.} A finite relational signature, a class of structures closed
under isomorphism, a subclass to be defined, and a logic with a complexity
bound. \emph{Conclusions.} Either a defining sentence within the bound, or a
proof that no such sentence exists, obtained from a winning strategy of the
second player in the appropriate game.

The presentation of first-order games and of the finite-variable logics
follows Kolaitis's manuscript \emph{On the Expressive Power of Logics on
Finite Models}, dated 1 August 2003 \cite{Kolaitis2003ExpressivePowerNotes}.
His chapter in the 2007 volume \emph{Finite Model Theory and Its
Applications} carries the same title; every locator given here refers to
the 2003 manuscript, and nothing is attributed to the printed chapter.
Locators for Gr\"adel's chapter in that volume refer to the author-hosted
copy of the chapter, whose running page numbers are those of the book
\cite{Gradel2007FMTDescriptiveChapter}. The manuscript outlines several of the
proofs or leaves them as hints; all proofs below are written out in full in
the annex. The basic Ehrenfeucht--Fra\"iss\'e game and quantifier rank are
those of \cref{sec:p03-ef-games}; the infinitary language and its ordinal
back-and-forth relations are those of \cref{sec:p03-w03-p03-039-marker}.

\section{Queries, definability and two measures of definitional complexity}
\label{sec:p03-w03-p03-055-measures}

\begin{hbconvention}{def:p03-w03-p03-055-setting}{Signatures, structures, assignments and partial isomorphisms}
Throughout, \(\tau\) is a finite set of relation symbols, each of positive
arity; there are no constant or function symbols. Structures are nonempty;
the universe of \(\mathcal A\) is written \(A\). The individual variables
are \(x_1,x_2,\ldots\), and \(X_k\defeq\{x_1,\ldots,x_k\}\). An
\emph{assignment} into \(\mathcal A\) is a map \(s\colon D\to A\) defined
on a finite set \(D\) of variables, and \(s[x\mapsto a]\) agrees with
\(s\) except that it sends \(x\) to \(a\) (its domain is
\(D\cup\{x\}\)). For assignments \(s\colon D\to A\) and
\(t\colon D\to B\) with the same domain put
\[
 \langle s,t\rangle\defeq\{(s(x),t(x)):x\in D\}\subseteq A\times B,
\]
and for tuples \(\bar a=(a_1,\ldots,a_n)\in A^n\),
\(\bar b\in B^n\) put \(\langle\bar a,\bar b\rangle\defeq
\{(a_i,b_i):1\le i\le n\}\); we write \(\mathcal A\entails\varphi[\bar a]\)
for satisfaction under \(x_i\mapsto a_i\).

A finite relation \(p\subseteq A\times B\) is a \emph{partial
isomorphism} if, for all \((a,b),(a',b')\in p\), \(a=a'\) exactly when
\(b=b'\), and for every \(R\in\tau\) of arity \(r\) and all
\((a_1,b_1),\ldots,(a_r,b_r)\in p\),
\(R^{\mathcal A}(a_1,\ldots,a_r)\) holds exactly when
\(R^{\mathcal B}(b_1,\ldots,b_r)\) holds. Equivalently, \(p\) is the graph
of an injective map from a finite subset of \(A\) onto a finite subset of
\(B\) that preserves and reflects every relation. Every subset of a partial
isomorphism is a partial isomorphism. This is the winning condition of
\cref{def:p03-ef-game}, and it agrees with Kolaitis's use of partial
isomorphisms as isomorphisms between generated substructures
\cite[p.~4 and Def.~3.1, p.~9]{Kolaitis2003ExpressivePowerNotes} in the
absence of constants.
\end{hbconvention}

\begin{hbdefinition}{def:p03-w03-p03-055-query}{Classes, Boolean queries and \(L\)-definability}
A \emph{class} of \(\tau\)-structures is a collection closed under
isomorphism. A \emph{Boolean query} on a class \(\mathcal C\) is an
isomorphism-invariant map \(Q\colon\mathcal C\to\{0,1\}\); it is identified
with the subclass \(\{\mathcal A\in\mathcal C:Q(\mathcal A)=1\}\)
\cite[Def.~2.1, p.~4]{Kolaitis2003ExpressivePowerNotes}. For a logic \(L\),
a subclass \(\mathcal K\subseteq\mathcal C\) is \emph{\(L\)-definable on
\(\mathcal C\)} if there is an \(L\)-sentence \(\sigma\) with
\(\mathcal A\in\mathcal K\iff\mathcal A\entails\sigma\) for every
\(\mathcal A\in\mathcal C\)
\cite[Def.~2.3, p.~5]{Kolaitis2003ExpressivePowerNotes}. One sentence must
work uniformly on all of \(\mathcal C\). Three choices of \(\mathcal C\) are
the important ones: all \(\tau\)-structures, a single structure up to
isomorphism, and the class \(\mathcal F\) of all finite \(\tau\)-structures.
Classifying a concept by definitional complexity means locating the least
logic, or the least bound inside a logic, for which it is definable on the
class in question.
\end{hbdefinition}

\begin{hbdefinition}{def:p03-w03-p03-055-logics}{Quantifier rank, finite-variable logics and the induced equivalences}
First-order formulas are built from \(\bottom\), atomic formulas
\(x=y\) and \(R(y_1,\ldots,y_r)\), the connectives \(\neg,\wedge,\vee\) and
the quantifiers \(\exists,\forall\). Quantifier rank \(\operatorname{qr}\) is
as in \cref{def:p03-quantifier-rank}
\cite[Def.~3.5, pp.~10--11]{Kolaitis2003ExpressivePowerNotes}.
\begin{enumerate}
\item \(\mathrm{FO}^k\) is the set of first-order formulas all of whose
      variables, free and bound, lie in \(X_k\).
\item \(\mathcal L^k_{\infty\omega}\) is the set of formulas of
      \(\mathcal L_{\infty,\omega}\) (\cref{def:p03-w03-p03-039-language})
      all of whose variables lie in \(X_k\); its formulas may use
      \(\bigwedge\Phi\) and \(\bigvee\Phi\) for arbitrary sets \(\Phi\) of
      such formulas. \(\mathcal L^\omega_{\infty\omega}\defeq
      \bigcup_{k\ge1}\mathcal L^k_{\infty\omega}\)
      \cite[Def.~7.1, pp.~52--53]{Kolaitis2003ExpressivePowerNotes}.
\item \(\mathcal A\equiv_m\mathcal B\) means that \(\mathcal A\) and
      \(\mathcal B\) satisfy the same first-order sentences of quantifier
      rank at most \(m\); \((\mathcal A,\bar a)\equiv_m(\mathcal B,\bar b)\)
      means the same for formulas with free variables among
      \(x_1,\ldots,x_n\), evaluated at \(\bar a\) and \(\bar b\)
      \cite[Def.~3.6, p.~11]{Kolaitis2003ExpressivePowerNotes}.
\item For assignments \(s,t\) with a common domain \(D\subseteq X_k\),
      \((\mathcal A,s)\equiv^k_{\infty\omega}(\mathcal B,t)\) means that
      \(\mathcal A\entails\varphi[s]\iff\mathcal B\entails\varphi[t]\) for
      every \(\mathcal L^k_{\infty\omega}\)-formula \(\varphi\) whose free
      variables lie in \(D\); \((\mathcal A,s)\equiv^k(\mathcal B,t)\) is the
      same with \(\mathrm{FO}^k\) in place of \(\mathcal L^k_{\infty\omega}\).
      For \(D=\emptyset\) these are written
      \(\mathcal A\equiv^k_{\infty\omega}\mathcal B\) and
      \(\mathcal A\equiv^k\mathcal B\)
      \cite[Def.~7.8, p.~57]{Kolaitis2003ExpressivePowerNotes}.
\end{enumerate}
A variable may be requantified any number of times, so an
\(\mathrm{FO}^k\)-formula can have arbitrarily large quantifier rank
\cite[p.~53]{Kolaitis2003ExpressivePowerNotes}. The truth of a formula of
either logic under an assignment depends only on the values of its free
variables.
\index{finite-variable logic}\index{quantifier rank!hierarchy}
\end{hbdefinition}

The two measures are related in one direction only.

\begin{hblemma}{thm:p03-w03-p03-055-rank-variables}{proved=annex:p03-w03-p03-055-rank-variables}{Rank bounds the number of variables}
Let \(\varphi\) be a first-order formula with free variables among
\(x_1,\ldots,x_n\) and \(\operatorname{qr}(\varphi)\le m\). Then
\(\varphi\) is logically equivalent to a formula of the same quantifier
rank whose variables, free and bound, lie in \(X_{n+m}\). In particular
every first-order sentence of quantifier rank at most \(m\ge1\) is
equivalent to an \(\mathrm{FO}^m\)-sentence.
\end{hblemma}

Gr\"adel states the same economy in terms of the \emph{width} of a formula,
the largest number of free variables of a subformula, and uses width and
quantifier rank to bound the space and time of first-order model checking
\cite[Lemma~3.1.4 and Thm.~3.1.5, p.~133]{Gradel2007FMTDescriptiveChapter}.
Knowing that a concept has a low-rank or few-variable definition therefore
also bounds the cost of evaluating it.

\section{Quantifier rank: definability and the game method}
\label{sec:p03-w03-p03-055-ef}

The Ehrenfeucht--Fra\"iss\'e theorem for positions with parameters is
proved in \cref{thm:p03-w03-p03-038-pointed-ef}: for a finite relational
signature, Duplicator wins the \(m\)-round game on \((\mathcal A,\bar a)\)
and \((\mathcal B,\bar b)\) exactly when the two pointed structures have the
same rank-\(m\) record, exactly when
\((\mathcal A,\bar a)\equiv_m(\mathcal B,\bar b)\). The records range over a
finite set \(\mathrm{Rec}_{n,m}\)
(\cref{thm:p03-w03-p03-038-finite-records}), and each record \(t\) is
defined by a characteristic formula \(\chi_t\) of rank at most \(m\)
(\cref{def:p03-w03-p03-038-characteristic,thm:p03-w03-p03-038-characteristic}).
The statement and its attribution to Fra\"iss\'e and Ehrenfeucht are
Kolaitis's Theorem~3.8 \cite[p.~11]{Kolaitis2003ExpressivePowerNotes}, and
the characteristic formulas correspond to his \((m,r)\)-types, which are
indexed by moves already made rather than by moves remaining
\cite[Def.~3.11 and Lemma~3.12, pp.~12--13]{Kolaitis2003ExpressivePowerNotes}.
The proved form of the game criterion is
\cref{thm:p03-w03-p03-038-criterion}. What the classification of concepts
by rank needs in addition is the converse: closure under \(\equiv_m\)
\emph{suffices} for definability by a sentence of rank \(m\), so the game
method is complete as well as sound.

\begin{hbtheorem}{thm:p03-w03-p03-055-rank-definability}{proved=annex:p03-w03-p03-055-rank-definability}{Definability by rank and completeness of the game method}
Let \(\mathcal C\) be a class of \(\tau\)-structures, finite or not, and
\(\mathcal K\subseteq\mathcal C\) a subclass closed under isomorphism; let
\(m\ge0\).
\begin{enumerate}
\item \(\mathcal K\) is definable on \(\mathcal C\) by a first-order
      sentence of quantifier rank at most \(m\) if and only if, for all
      \(\mathcal A,\mathcal B\in\mathcal C\) with
      \(\mathcal A\equiv_m\mathcal B\), \(\mathcal A\in\mathcal K\) implies
      \(\mathcal B\in\mathcal K\). The defining sentence can be taken to be a
      disjunction of characteristic sentences \(\chi_t\),
      \(t\in\mathrm{Rec}_{0,m}\).
\item \(\mathcal K\) is not first-order definable on \(\mathcal C\) if and
      only if for every \(m\) there are \(\mathcal A_m\in\mathcal K\) and
      \(\mathcal B_m\in\mathcal C\setminus\mathcal K\) such that Duplicator
      wins the \(m\)-round game on \(\mathcal A_m\) and \(\mathcal B_m\).
\end{enumerate}
\end{hbtheorem}

These are Kolaitis's Theorem~3.15 and Method~3.16
\cite[pp.~14--15]{Kolaitis2003ExpressivePowerNotes}. The ``if'' half of the
second item, for the class of finite structures, is
\cref{thm:p03-ef-inexpressibility} and
\cref{thm:p03-w03-p03-038-criterion}; the ``only if'' half says that every
first-order undefinability fact on \(\mathcal C\) has a game proof inside
\(\mathcal C\).

\begin{hbexample}{def:p03-w03-p03-055-ef-strategy-reading}{Reading a distinguishing sentence as a strategy}
A distinguishing sentence is a recipe for Spoiler. If
\(\sigma=\exists x_1\exists x_2\exists x_3\bigwedge_{i<j}
(\neg x_i=x_j\wedge\neg E(x_i,x_j))\) holds in a graph \(\mathcal B\) and
fails in \(\mathcal A\), Spoiler plays three pairwise distinct, pairwise
nonadjacent vertices of \(\mathcal B\). Because \(\mathcal A\) has no such
triple, two of Duplicator's three answers in \(\mathcal A\) are equal or
adjacent, and the selected pairing is not a partial isomorphism: Spoiler
wins the \(3\)-round game, as \cref{thm:p03-w03-p03-038-pointed-ef} predicts
from \(\operatorname{qr}(\sigma)=3\). Kolaitis discusses a pair of graphs
separated in this way at three rounds but not at two
\cite[Example~3.3 and Example~3.9, pp.~9--11]{Kolaitis2003ExpressivePowerNotes}.
\end{hbexample}

\section{Finite-variable logics and pebble games}
\label{sec:p03-w03-p03-055-pebble}

\begin{hbdefinition}{def:p03-w03-p03-055-pebble}{The \(k\)-pebble game and winning families}
Let \(k\ge1\). In the \emph{\(k\)-pebble game} on \(\mathcal A\) and
\(\mathcal B\) each player has \(k\) pebbles labelled \(1,\ldots,k\). In
each move Spoiler, in a structure of his choice, either places an unused
pebble on an element or removes a pebble; Duplicator places or removes the
pebble with the same label in the other structure. Spoiler wins as soon as
the pairs of elements carrying equally labelled pebbles fail to form a
partial isomorphism; Duplicator wins if this never happens, however long
the play continues \cite[Def.~7.3, p.~55]{Kolaitis2003ExpressivePowerNotes}.
Following Kolaitis's precise form
\cite[Def.~7.4, p.~55]{Kolaitis2003ExpressivePowerNotes}, a \emph{winning
family} for Duplicator is a nonempty set \(\mathcal I\) of partial
isomorphisms from \(\mathcal A\) to \(\mathcal B\) such that
\begin{enumerate}
\item \(|f|\le k\) for every \(f\in\mathcal I\);
\item if \(g\in\mathcal I\) and \(f\subseteq g\) then \(f\in\mathcal I\);
\item (forth up to \(k\)) if \(f\in\mathcal I\), \(|f|<k\) and
      \(a\in A\), there is \(g\in\mathcal I\) with \(f\subseteq g\) and
      \(a\in\dom g\);
\item (back up to \(k\)) if \(f\in\mathcal I\), \(|f|<k\) and
      \(b\in B\), there is \(g\in\mathcal I\) with \(f\subseteq g\) and
      \(b\) in the range of \(g\).
\end{enumerate}
Duplicator \emph{wins the \(k\)-pebble game} on \(\mathcal A\) and
\(\mathcal B\) when a winning family exists (it then contains \(\emptyset\)
by the second clause), and on \((\mathcal A,s)\) and \((\mathcal B,t)\),
for assignments with a common domain \(D\subseteq X_k\), when some winning
family contains \(\langle s,t\rangle\). Clause~2 answers the removal of a
pebble, clauses~3 and~4 answer its placement.
\end{hbdefinition}

\begin{hbtheorem}{thm:p03-w03-p03-055-pebble}{proved=annex:p03-w03-p03-055-pebble}{Pebble games characterize \(k\)-variable equivalence}
Let \(k\ge1\), let \(\mathcal A,\mathcal B\) be \(\tau\)-structures, and
let \(s\colon D\to A\), \(t\colon D\to B\) with \(D\subseteq X_k\).
The following are equivalent:
\begin{enumerate}
\item Duplicator wins the \(k\)-pebble game on \((\mathcal A,s)\) and
      \((\mathcal B,t)\);
\item \((\mathcal A,s)\equiv^k_{\infty\omega}(\mathcal B,t)\).
\end{enumerate}
If \(\mathcal A\) and \(\mathcal B\) are finite, both are also equivalent to
\begin{enumerate}\setcounter{enumi}{2}
\item \((\mathcal A,s)\equiv^k(\mathcal B,t)\).
\end{enumerate}
\end{hbtheorem}

This is Kolaitis's Theorem~7.9, attributed to Barwise and to Immerman
\cite[pp.~58--59]{Kolaitis2003ExpressivePowerNotes}; the manuscript gives an
outline and refers to Kolaitis and Vardi for the details. The annex supplies
the complete argument. The finiteness hypothesis is used at exactly one
point: it makes the conjunctions in the back-and-forth step finite, so that
they stay inside \(\mathrm{FO}^k\).

\begin{hbtheorem}{thm:p03-w03-p03-055-k-definability}{proved=annex:p03-w03-p03-055-k-definability}{Definability in \(k\)-variable infinitary logic on finite structures}
Let \(\mathcal C\) be a class of finite \(\tau\)-structures,
\(\mathcal K\subseteq\mathcal C\) closed under isomorphism, and \(k\ge1\).
Then \(\mathcal K\) is \(\mathcal L^k_{\infty\omega}\)-definable on
\(\mathcal C\) if and only if, for all \(\mathcal A,\mathcal B\in\mathcal C\)
with \(\mathcal A\equiv^k\mathcal B\) (equivalently, with Duplicator
winning the \(k\)-pebble game on them), \(\mathcal A\in\mathcal K\) implies
\(\mathcal B\in\mathcal K\). Hence \(\mathcal K\) is
\(\mathcal L^\omega_{\infty\omega}\)-definable on \(\mathcal C\) if and only
if this holds for some \(k\).
\end{hbtheorem}

This is Kolaitis's Corollary~7.10
\cite[p.~59]{Kolaitis2003ExpressivePowerNotes}. The bound on variables is
what gives the theorem content.

\begin{hblemma}{thm:p03-w03-p03-055-unbounded}{proved=annex:p03-w03-p03-055-unbounded}{Without a variable bound every finite-structure class is definable}
Every class \(\mathcal K\) of finite \(\tau\)-structures is definable on the
class \(\mathcal F\) of all finite \(\tau\)-structures by a sentence of
\(\mathcal L_{\infty,\omega}\) that is a countable disjunction of
first-order sentences.
\end{hblemma}

Kolaitis records this observation at the opening of his Section~7
\cite[p.~52]{Kolaitis2003ExpressivePowerNotes}. The defining disjunction
uses unboundedly many variables, so the lemma places no concept in any
\(\mathcal L^k_{\infty\omega}\).

\begin{hblemma}{thm:p03-w03-p03-055-pebble-basic}{proved=annex:p03-w03-p03-055-pebble-basic}{Monotonicity in the number of pebbles}
Let \(1\le k\le k'\). If Duplicator wins the \(k'\)-pebble game on
\(\mathcal A\) and \(\mathcal B\), she wins the \(k\)-pebble game on them;
equivalently, if Spoiler wins the \(k\)-pebble game, he wins the
\(k'\)-pebble game. For each \(k\), the relation ``Duplicator wins the
\(k\)-pebble game on \(\mathcal A\) and \(\mathcal B\)'' is an equivalence
relation on the class of all \(\tau\)-structures.
\end{hblemma}

Kolaitis records both facts after his definition of the game
\cite[p.~56]{Kolaitis2003ExpressivePowerNotes}.

\section{Rounds and pebbles together}
\label{sec:p03-w03-p03-055-bounded}

The two measures combine. A formula of \(\mathrm{FO}^k\) of quantifier rank
at most \(m\) is bounded in both resources at once, and the game that
matches it limits both the pebbles and the number of rounds. The resulting
two-parameter family contains the Ehrenfeucht--Fra\"iss\'e game as the case
in which there are at least as many pebbles as rounds, and the unbounded
pebble game of \cref{def:p03-w03-p03-055-pebble} as the limit in which the
number of rounds is unrestricted.

\begin{hbdefinition}{def:p03-w03-p03-055-bounded-pebble}{The \(m\)-round \(k\)-pebble game and \(\equiv^k_m\)}
Let \(k\ge1\) and \(m\ge0\). A \emph{position} on \(\mathcal A\) and
\(\mathcal B\) is a pair \((s,t)\) of assignments \(s\colon D\to A\),
\(t\colon D\to B\) with a common domain \(D\subseteq X_k\); the variable
\(x_i\in D\) records where pebble \(i\) lies in each structure. In one
\emph{round} Spoiler chooses a label \(i\in\{1,\ldots,k\}\) and an element
\(a\in A\) (or \(b\in B\)); Duplicator answers with an element \(b\in B\)
(or \(a\in A\)), and the new position is
\((s[x_i\mapsto a],t[x_i\mapsto b])\). A pebble already on the board is
thus lifted and replaced in a single round. Duplicator \emph{wins
\(G^k_m(\mathcal A,s;\mathcal B,t)\)}, the \(m\)-round \(k\)-pebble game
from \((s,t)\), as defined by recursion on \(m\):
\begin{enumerate}
\item she wins \(G^k_0(\mathcal A,s;\mathcal B,t)\) if
      \(\langle s,t\rangle\) is a partial isomorphism;
\item she wins \(G^k_{m+1}(\mathcal A,s;\mathcal B,t)\) if
      \(\langle s,t\rangle\) is a partial isomorphism, for every
      \(i\in\{1,\ldots,k\}\) and \(a\in A\) there is \(b\in B\) such that
      she wins \(G^k_m(\mathcal A,s[x_i\mapsto a];\mathcal B,t[x_i\mapsto b])\),
      and for every \(i\) and \(b\in B\) there is \(a\in A\) with the same
      property.
\end{enumerate}
This recursion says that Duplicator has a strategy that keeps the pebbled
pairs a partial isomorphism at the start and after each of the \(m\)
rounds. For assignments with common domain \(D\subseteq X_k\),
\((\mathcal A,s)\equiv^k_m(\mathcal B,t)\) means that
\(\mathcal A\entails\varphi[s]\iff\mathcal B\entails\varphi[t]\) for every
\(\mathrm{FO}^k\)-formula \(\varphi\) of quantifier rank at most \(m\) whose
free variables lie in \(D\); for \(D=\emptyset\) we write
\(\mathcal A\equiv^k_m\mathcal B\). Two formulas are \emph{logically
equivalent} when they hold under the same assignments of all structures.
\index{pebble game!bounded}
\end{hbdefinition}

\begin{hblemma}{thm:p03-w03-p03-055-finitely-many-k}{proved=annex:p03-w03-p03-055-finitely-many-k}{Finitely many \(k\)-variable formulas of each rank}
Let \(\tau\) be finite, \(k\ge1\), \(m\ge0\) and \(D\subseteq X_k\). Up to
logical equivalence there are only finitely many \(\mathrm{FO}^k\)-formulas
of quantifier rank at most \(m\) with free variables in \(D\).
\end{hblemma}

\begin{hbtheorem}{thm:p03-w03-p03-055-bounded-pebble}{proved=annex:p03-w03-p03-055-bounded-pebble}{The \(m\)-round \(k\)-pebble game characterizes \(\equiv^k_m\)}
Let \(\tau\) be finite, \(k\ge1\), \(m\ge0\), let \(\mathcal A,\mathcal B\)
be \(\tau\)-structures, finite or infinite, and let \(s\colon D\to A\),
\(t\colon D\to B\) with \(D\subseteq X_k\). Then Duplicator wins
\(G^k_m(\mathcal A,s;\mathcal B,t)\) if and only if
\((\mathcal A,s)\equiv^k_m(\mathcal B,t)\).
\end{hbtheorem}

The implication from the game to agreement does not use the finiteness of
\(\tau\); the converse uses it through
\cref{thm:p03-w03-p03-055-finitely-many-k}, which keeps the conjunction that
Spoiler's move expresses finite.

\begin{hbcorollary}{thm:p03-w03-p03-055-bounded-consequences}{proved=annex:p03-w03-p03-055-bounded-consequences}{Consequences of the two-parameter game}
Let \(\tau\) be finite and \(k\ge1\).
\begin{enumerate}
\item For all \(\tau\)-structures and assignments with common domain
      \(D\subseteq X_k\), \((\mathcal A,s)\equiv^k(\mathcal B,t)\) if and
      only if Duplicator wins \(G^k_m(\mathcal A,s;\mathcal B,t)\) for every
      \(m\). If \(\mathcal A\) and \(\mathcal B\) are finite, this holds if
      and only if Duplicator wins the \(k\)-pebble game on
      \((\mathcal A,s)\) and \((\mathcal B,t)\).
\item If \(k\ge m\), then \(\mathcal A\equiv^k_m\mathcal B\) if and only if
      \(\mathcal A\equiv_m\mathcal B\), if and only if Duplicator wins the
      \(m\)-round Ehrenfeucht--Fra\"iss\'e game on \(\mathcal A\) and
      \(\mathcal B\).
\item Let \(\mathcal C\) be a class of \(\tau\)-structures and
      \(\mathcal K\subseteq\mathcal C\) closed under isomorphism. Then
      \(\mathcal K\) is definable on \(\mathcal C\) by an
      \(\mathrm{FO}^k\)-sentence of quantifier rank at most \(m\) if and only
      if, for all \(\mathcal A,\mathcal B\in\mathcal C\) with
      \(\mathcal A\equiv^k_m\mathcal B\), \(\mathcal A\in\mathcal K\)
      implies \(\mathcal B\in\mathcal K\). Consequently \(\mathcal K\) is
      \(\mathrm{FO}^k\)-definable on \(\mathcal C\) if and only if this
      holds for some \(m\), and it is not \(\mathrm{FO}^k\)-definable on
      \(\mathcal C\) if and only if for every \(m\) there are
      \(\mathcal A_m\in\mathcal K\) and
      \(\mathcal B_m\in\mathcal C\setminus\mathcal K\) such that Duplicator
      wins \(G^k_m(\mathcal A_m,\emptyset;\mathcal B_m,\emptyset)\).
\end{enumerate}
\end{hbcorollary}

The second item locates the rank hierarchy inside the two-parameter one: up
to logical equivalence the sentences of rank at most \(m\) are already
\(\mathrm{FO}^m\)-sentences (\cref{thm:p03-w03-p03-055-rank-variables}), so
pebbles beyond the number of rounds add nothing. The first item locates the
variable hierarchy: on finite structures the unbounded \(k\)-pebble game is
the conjunction of all its round-bounded versions. For infinite structures
the round-bounded games still capture \(\equiv^k\), while the unbounded game
is tied to \(\equiv^k_{\infty\omega}\) (\cref{thm:p03-w03-p03-055-pebble});
no comparison of the two relations on infinite structures is proved here.

\section{The two hierarchies and worked cases}
\label{sec:p03-w03-p03-055-hierarchies}

Let \(K_n\) be the complete graph on \(n\) vertices (the \(n\)-clique),
\(L_n=(\{1,\ldots,n\},<)\) the \(n\)-element linear order, and, for
\(n\ge3\), \(C_n\) the cycle with vertex set \(\mathbb Z/n\mathbb Z\) and
edges \(\{i,i+1\}\); graphs are structures for one binary symbol \(E\)
interpreted symmetrically and irreflexively. Put
\[
 \sigma_{\ge n}\defeq\exists x_1\cdots\exists x_n
 \bigwedge_{1\le i<j\le n}\neg x_i=x_j ,
\]
a sentence of \(\mathrm{FO}^n\) of quantifier rank \(n\) that says there are
at least \(n\) elements.

\begin{hbtheorem}{thm:p03-w03-p03-055-cliques}{proved=annex:p03-w03-p03-055-cliques}{Cliques and the strictness of both hierarchies}
Let \(k\ge1\).
\begin{enumerate}
\item For \(n,n'\ge k\) the set of all injective partial maps
      \(f\) from \(K_n\) to \(K_{n'}\) with \(|f|\le k\) is a winning family
      in the \(k\)-pebble game; hence
      \(K_n\equiv^k_{\infty\omega}K_{n'}\) and \(K_n\equiv_kK_{n'}\).
\item Spoiler wins the \((k+1)\)-pebble game and the \((k+1)\)-round game
      on \(K_k\) and \(K_{k+1}\).
\item (Variable hierarchy.) The class of graphs with at least \(k+1\)
      vertices is defined on finite graphs by the \(\mathrm{FO}^{k+1}\)-sentence
      \(\sigma_{\ge k+1}\) but is not \(\mathcal L^k_{\infty\omega}\)-definable
      on finite graphs, nor even on finite cliques.
\item (Rank hierarchy.) The same class is not definable on finite cliques
      by any first-order sentence of quantifier rank at most \(k\).
\end{enumerate}
\end{hbtheorem}

The first two items are Kolaitis's Example~7.5 together with his remark on
the corresponding round counts
\cite[p.~56]{Kolaitis2003ExpressivePowerNotes}.

\begin{hbtheorem}{thm:p03-w03-p03-055-orders}{proved=annex:p03-w03-p03-055-orders}{Two variables on finite linear orders}
Define \(\mathrm{FO}^2\)-formulas over \(\{<\}\) by
\(\alpha_1(x_1)\defeq x_1=x_1\), \(\beta_1(x_2)\defeq x_2=x_2\) and
\[
 \alpha_{j+1}(x_1)\defeq\exists x_2\bigl(x_2<x_1\wedge\beta_j(x_2)\bigr),
 \qquad
 \beta_{j+1}(x_2)\defeq\exists x_1\bigl(x_1<x_2\wedge\alpha_j(x_1)\bigr).
\]
\begin{enumerate}
\item In every linear order, \(\alpha_j\) and \(\beta_j\) hold of an element
      exactly when at least \(j-1\) elements lie below it; so
      \(\exists x_1\,\alpha_n(x_1)\) is an \(\mathrm{FO}^2\)-sentence of
      quantifier rank \(n\) equivalent on linear orders to
      \(\sigma_{\ge n}\).
\item For \(m<n\), Spoiler wins the \(2\)-pebble game on \(L_m\) and
      \(L_n\).
\item For every set \(S\) of positive integers, the class of finite linear
      orders whose size lies in \(S\) is
      \(\mathcal L^2_{\infty\omega}\)-definable on finite linear orders. In
      particular even cardinality is, although it is not first-order
      definable on finite linear orders (\cref{def:p03-ef-parity}).
\end{enumerate}
\end{hbtheorem}

The \(\mathrm{FO}^2\)-sentence for four elements, the definability of every
set of cardinalities and Spoiler's two-pebble strategy are Kolaitis's
\cite[p.~53 and Example~7.6, p.~56]{Kolaitis2003ExpressivePowerNotes}. The
contrast with the parity example is the point: for a fixed number of rounds
Duplicator wins the Ehrenfeucht--Fra\"iss\'e game on all sufficiently long
finite linear orders, while two pebbles, reused, separate every pair of
distinct lengths. Quantifier rank and number of variables measure different
resources, and neither hierarchy refines the other.

Connectivity shows rank and variables parting company.
\Cref{thm:p03-w03-p03-038-connectivity} proves, with the distance invariant
and response lemma of \cref{def:p03-w03-p03-038-cycles} and
\cref{thm:p03-w03-p03-038-cycle-response}, that Duplicator wins the
\(k\)-round game on \(C_{2m}\) and \(C_m\sqcup C_m\) whenever
\(m\ge m_k=(k+1)2^{k+2}\); this is the strategy recorded in
\cref{def:p03-ef-connectivity}. Kolaitis reaches the same pairs of graphs
through a Hanf-type counting condition whose proof he gives as a hint
\cite[Thm.~3.26 and Props.~3.28--3.29, pp.~19--20]{Kolaitis2003ExpressivePowerNotes}.

\begin{hbcorollary}{thm:p03-w03-p03-055-connectivity}{proved=annex:p03-w03-p03-055-connectivity}{Connectivity and 2-colorability}
\leavevmode
\begin{enumerate}
\item 2-colorability is not first-order definable on the class of
      finite graphs; the same holds for connectivity
      (\cref{thm:p03-w03-p03-038-connectivity}).
\item Connectivity is \(\mathcal L^3_{\infty\omega}\)-definable on all
      graphs: with \(\pi_1(x_1,x_2)\defeq E(x_1,x_2)\) and
      \[
       \pi_{n+1}(x_1,x_2)\defeq\exists x_3\bigl(E(x_1,x_3)\wedge
       \exists x_1(x_3=x_1\wedge\pi_n(x_1,x_2))\bigr),
      \]
      \(\pi_n\) expresses the existence of a walk of length exactly \(n\),
      and a graph is connected exactly when it satisfies
      \(\forall x_1\forall x_2\bigl(x_1=x_2\vee\bigvee_{n\ge1}
      \pi_n(x_1,x_2)\bigr)\).
\item For every \(m\ge3\), Spoiler wins the \(3\)-pebble game on
      \(C_{2m}\) and \(C_m\sqcup C_m\).
\end{enumerate}
\end{hbcorollary}

The three-variable path formulas are Kolaitis's
\cite[pp.~53--54]{Kolaitis2003ExpressivePowerNotes}; his sentence omits the
disjunct \(x_1=x_2\) and so is false on the one-vertex graph, which is
connected. His Example~7.7 makes the same three-pebble point for directed
cycles \cite[pp.~56--57]{Kolaitis2003ExpressivePowerNotes}.

\begin{hblemma}{thm:p03-w03-p03-055-bipartite}{proved=annex:p03-w03-p03-055-bipartite}{2-colorability in three-variable infinitary logic}
A graph is 2-colorable if and only if it has no closed walk of odd length.
Consequently 2-colorability is defined on the class of all graphs by the
\(\mathcal L^3_{\infty\omega}\)-sentence
\[
 \neg\exists x_1\exists x_2\Bigl(x_1=x_2\wedge
 \bigvee_{n\ge1,\ n\text{ odd}}\pi_n(x_1,x_2)\Bigr),
\]
with \(\pi_n\) as in \cref{thm:p03-w03-p03-055-connectivity}, and for every
odd \(m\ge3\) Spoiler wins the \(3\)-pebble game on \(C_{2m}\) and
\(C_m\sqcup C_m\).
\end{hblemma}

The proof of the characterization uses the axiom of choice to pick one
vertex in each component of an infinite graph; for finite graphs no choice
is needed.

\begin{hbmethod}{def:p03-w03-p03-055-classify}{Classifying a concept of finite structures}
Given a finite relational signature \(\tau\), a class \(\mathcal C\) of
\(\tau\)-structures closed under isomorphism, and a subclass
\(\mathcal K\) closed under isomorphism:
\begin{enumerate}
\item Fix \(\mathcal C\) first. Every upper and lower bound below is
      relative to it; a result for all finite structures does not transfer
      to a smaller class, nor conversely.
\item Upper bound. Write a defining sentence and read off its quantifier
      rank and its set of variables. Rename bound variables with
      \cref{thm:p03-w03-p03-055-rank-variables} when only the rank matters,
      and reuse variables, as in \cref{thm:p03-w03-p03-055-orders} and
      \cref{thm:p03-w03-p03-055-connectivity}, when only their number
      matters.
\item Lower bound for rank \(m\). Exhibit \(\mathcal A\in\mathcal K\) and
      \(\mathcal B\in\mathcal C\setminus\mathcal K\) with a Duplicator
      strategy for the \(m\)-round game
      (\cref{thm:p03-w03-p03-055-rank-definability}). For non-definability
      in first-order logic such a pair is needed for every \(m\).
\item Lower bound for \(k\) variables. On finite structures exhibit a pair
      as in the previous step with a winning family for the \(k\)-pebble
      game (\cref{thm:p03-w03-p03-055-k-definability}); this excludes
      \(\mathcal L^k_{\infty\omega}\) and hence \(\mathrm{FO}^k\). To
      exclude only \(\mathrm{FO}^k\), a family of pairs on which Duplicator
      wins \(G^k_m\) for every \(m\) suffices
      (\cref{thm:p03-w03-p03-055-bounded-consequences}).
\item Side conditions. The signature must be finite for the completeness
      halves (the converse directions) of
      \cref{thm:p03-w03-p03-055-rank-definability,thm:p03-w03-p03-055-bounded-consequences};
      the soundness halves need only that the strategies are genuine. The
      identification of the \(k\)-pebble game with \(\equiv^k\) requires
      finite structures.
\end{enumerate}
\end{hbmethod}

\begin{hbexample}{def:p03-w03-p03-055-classification}{A classification of six concepts}
The results of this chapter place the following concepts. The class
\(\mathcal C\) is stated in each case.
\begin{center}
\begin{tabular}{p{0.3\linewidth}p{0.6\linewidth}}
\hline
Concept and class & Definability \\
\hline
At least \(n\) elements, finite graphs, \(n\ge2\) &
\(\mathrm{FO}^n\), rank \(n\); not \(\mathcal L^{n-1}_{\infty\omega}\); not
rank \(n-1\) (\cref{thm:p03-w03-p03-055-cliques}) \\
At least \(n\) elements, finite linear orders &
\(\mathrm{FO}^2\), rank \(n\) (\cref{thm:p03-w03-p03-055-orders}) \\
Size in a set \(S\), finite linear orders &
\(\mathcal L^2_{\infty\omega}\) for every \(S\)
(\cref{thm:p03-w03-p03-055-orders}) \\
Even size, finite linear orders &
\(\mathcal L^2_{\infty\omega}\); not first-order
(\cref{thm:p03-w03-p03-055-orders}, \cref{def:p03-ef-parity}) \\
Connectivity, finite graphs &
\(\mathcal L^3_{\infty\omega}\); not first-order
(\cref{thm:p03-w03-p03-055-connectivity}) \\
2-colorability, finite graphs &
\(\mathcal L^3_{\infty\omega}\); not first-order
(\cref{thm:p03-w03-p03-055-connectivity,thm:p03-w03-p03-055-bipartite}) \\
\hline
\end{tabular}
\end{center}
For the first row with \(n=k+1\) the upper bound is the sentence
\(\sigma_{\ge k+1}\), and both lower bounds come from the single pair
\(K_k,K_{k+1}\). The rows for connectivity and 2-colorability use the same
pairs \(C_{2m},C_m\sqcup C_m\) in opposite roles: Duplicator wins the
\(k\)-round game when \(m\) is large (and odd, for 2-colorability), while
Spoiler wins the \(3\)-pebble game for every \(m\ge3\) (every odd
\(m\ge3\), for 2-colorability).
\end{hbexample}

\begin{hbwarning}{def:p03-w03-p03-055-limits}{What the game methods do and do not give}
\begin{enumerate}
\item A nondefinability proof by games concerns a stated class: the pairs
      \(\mathcal A_m,\mathcal B_m\) must lie in \(\mathcal C\). Even
      cardinality is not first-order definable on finite linear orders, yet
      it is \(\mathcal L^2_{\infty\omega}\)-definable there
      (\cref{thm:p03-w03-p03-055-orders}); the two statements concern
      different logics on the same class and do not conflict.
\item The equivalence of \(\equiv^k\) with the \(k\)-pebble game is proved
      here only for finite structures. For infinite structures the proof
      gives only the equivalence of the game with
      \(\equiv^k_{\infty\omega}\); no claim is made that \(\equiv^k\) and
      \(\equiv^k_{\infty\omega}\) agree there.
\item The finite-variable characterization does not bound quantifier rank,
      and the rank characterization does not bound variables:
      \cref{thm:p03-w03-p03-055-orders} has \(\mathrm{FO}^2\)-sentences of
      every rank, while \(\sigma_{\ge k+1}\) has rank \(k+1\) and needs
      \(k+1\) variables on cliques.
\item \Cref{thm:p03-w03-p03-055-unbounded} shows that on finite
      structures \(\mathcal L_{\infty,\omega}\) without a variable bound
      separates nothing: classification by infinitary logic is informative
      only through \(k\).
\end{enumerate}
\end{hbwarning}

\paragraph{Relationships.}
\Cref{thm:p03-w03-p03-055-rank-definability} completes the game method of
\cref{sec:p03-ef-games} and of \cref{sec:p03-w03-p03-038-comparison-games}
by its converse, and the transfinite back-and-forth relations used for
Karp's and Scott's theorems are in \cref{sec:p03-w03-p03-039-marker}. The pebble game is the
\(k\)-variable analogue of the partial-isomorphism systems of
\cref{def:p03-w03-p03-039-partial-iso}. Asymptotic laws for finite
structures, which also rest on extension and back-and-forth arguments, are
in \cref{sec:p03-finite-model-theory}. Gr\"adel's chapter places these
definability questions within descriptive complexity, where logics are
compared with complexity classes on finite structures
\cite[Section~3.1, pp.~125--138]{Gradel2007FMTDescriptiveChapter}.
